\section{Verificação e Ambiente}
\label{sec:verificacao-ambiente}

\subsection{Testes Automatizados com Pytest}
\label{subsec:testes-automatizados-pytest}



O Pytest foi empregado para verificar as regras de negócio do módulo, em especial as transições
permitidas da máquina de estados: criação, aprovação, rejeição e reversão \cite{pytest2024}.
Os testes também verificaram os efeitos dessas transições sobre a categoria, a pontuação e o
histórico da solicitação.

% TO-DO: Imagem de uma execução de teste.
\todo[inline]{TODO:Imagem de uma execução de teste.}

\subsection{Testes Unitários}
\label{subsec:testes-unitarios}
Teste unitário é aquele que verifica, de forma isolada, o comportamento de uma
unidade do sistema, como uma função, uma classe ou um serviço. Seu objetivo é
avaliar uma regra específica sem depender da comunicação com outros
componentes, o que torna mais precisa a identificação de falhas \cite{pytest2024}.

\subsection{Testes de Integração}
\label{subsec:testes-integracao}
Teste de integração, por sua vez, verifica a interação entre dois ou mais
componentes, avaliando se eles se comunicam e funcionam corretamente em
conjunto. Diferentemente do teste unitário, pode envolver recursos externos,
como banco de dados, API e ambiente de execução, para validar o fluxo completo
da aplicação.

\subsection{Ambiente Reprodutível com Docker}
\label{subsec:ambiente-reprodutivel-docker}

O Docker foi utilizado para reproduzir o ambiente necessário aos testes de integração, incluindo
a API e o PostgreSQL, de modo que a validação dos fluxos ocorresse sob uma configuração
consistente entre execuções \cite{docker2024}.
